\documentclass[12pt]{article}
\usepackage[utf8]{inputenc}
\usepackage{amsmath, amssymb}
\usepackage{geometry}
\geometry{margin=1in}

\title{The Specific-Factors Model}
\author{ECON 308 -- International Economic Relations}
\date{}

\begin{document}
\maketitle



Home produces Manufactures ($M$) and Food ($F$). Labor is mobile between sectors, while capital is specific to $M$ and land is specific to $F$.  

\[
Q_M = K^{1/2}L_M^{1/2}, 
\qquad 
Q_F = T^{1/2}L_F^{1/2},
\qquad 
L_M + L_F = L.
\]

Endowments:
\[
K=100, \quad T=64, \quad L=100.
\]

Prices: initially $P_M=2$, $P_F=1$. Later, $P_M$ rises to $3$ while $P_F$ remains $1$.  

\bigskip

\noindent 1. Derive the marginal products of labor in each sector.  

\bigskip

\noindent 2. Write the wage equalization condition and show the equation that determines the allocation of labor between the two sectors.  

\bigskip

\noindent 3. Solve for the equilibrium allocation of labor ($L_M$, $L_F$) and the common wage when $P_M=2$.  

\bigskip

\noindent 4. Recalculate the equilibrium labor allocation and wage when $P_M=3$.  

\bigskip

\noindent 5. Compare the results. Discuss what happens to $L_M$, $w$, and the real wages of labor in terms of both goods. Identify which groups (capital owners, landowners, or workers) gain or lose from the price change.

\end{document}
